\title{2017 年高考数学 全国II卷-理科}
\date{}
\maketitle

\section*{一、选择题}
本题共 12 小题，每小题 5 分，共 60 分。在每小题给出的四个选项中，只有一项是符合题目要求的。

\begin{question}[points=5]
$\frac{3+i}{1+i}=$（   ）
\begin{choices}
  \item $1+2i$
  \item $1-2i$
  \item $2+i$
  \item $2-i$
\end{choices}
\begin{solution}
\textbf{答案：}$D$

\textbf{分析：}分子和分母同时乘以分母的共轭复数，再利用虚数单位 $i$ 的幂运算性质，求出结果。

\textbf{详解：}$\frac{3+i}{1+i}=\frac{(3+i)(1-i)}{(1+i)(1-i)}=\frac{4-2i}{2}=2-i$，故选：$D$。
\end{solution}
\end{question}

\begin{question}[points=5]
设集合 $A=\{1,2,4\}$，$B=\{x|x^2-4x+m=0\}$．若 $A\cap B=\{1\}$，则 $B=$（   ）
\begin{choices}
  \item $\{1,-3\}$
  \item $\{1,0\}$
  \item $\{1,3\}$
  \item $\{1,5\}$
\end{choices}
\begin{solution}
\textbf{答案：}$C$

\textbf{分析：}由交集的定义可得 $1\in A$ 且 $1\in B$，代入二次方程，求得 $m$，再解二次方程可得集合 $B$。

\textbf{详解：}集合 $A=\{1,2,4\}$，$B=\{x|x^2-4x+m=0\}$．若 $A\cap B=\{1\}$，则 $1\in A$ 且 $1\in B$，可得 $1-4+m=0$，解得 $m=3$，即有 $B=\{x|x^2-4x+3=0\}=\{1,3\}$．故选：$C$。
\end{solution}
\end{question}

\begin{question}[points=5]
我国古代数学名著《算法统宗》中有如下问题：“远看巍巍塔七层，红光点点倍加增，共灯三百八十一，请问尖头几盏灯？”意思是：一座 $7$ 层塔共挂了 $381$ 盏灯，且相邻两层中的下一层灯数是上一层灯数的 $2$ 倍，则塔的顶层共有灯（   ）
\begin{choices}
  \item $1$ 盏
  \item $3$ 盏
  \item $5$ 盏
  \item $9$ 盏
\end{choices}
\begin{solution}
\textbf{答案：}$B$

\textbf{分析：}设塔顶的 $a_1$ 盏灯，由题意 $\{a_n\}$ 是公比为 $2$ 的等比数列，利用等比数列前 $n$ 项和公式列出方程，能求出结果。

\textbf{详解：}设塔顶的 $a_1$ 盏灯，由题意 $\{a_n\}$ 是公比为 $2$ 的等比数列，$\therefore S_7=\frac{a_1(1-2^7)}{1-2}=381$，解得 $a_1=3$．故选：$B$。
\end{solution}
\end{question}

\begin{question}[points=5]
如图，网格纸上小正方形的边长为 $1$，粗实线画出的是某几何体的三视图，该几何体由一平面将一圆柱截去一部分后所得，则该几何体的体积为（   ）

\begin{center}
  \includegraphics[width=.38\linewidth]{question-4-1.jpg}
\end{center}
\begin{choices}
  \item $90\pi$
  \item $63\pi$
  \item $42\pi$
  \item $36\pi$
\end{choices}
\begin{solution}
\textbf{答案：}$B$

\textbf{分析：}由三视图可得，直观图为一个完整的圆柱减去一个高为 $6$ 的圆柱的一半，即可求出几何体的体积。

\textbf{详解：}由三视图可得，直观图为一个完整的圆柱减去一个高为 $6$ 的圆柱的一半，$V=\pi\cdot3^2\times10-\frac{1}{2}\pi\cdot3^2\times6=63\pi$，故选：$B$。
\end{solution}
\end{question}

\begin{question}[points=5]
设 $x$，$y$ 满足约束条件 $\begin{cases}2x+3y-3\le 0\\2x-3y+3\ge 0\\y+3\ge 0\end{cases}$，则 $z=2x+y$ 的最小值是（   ）

\begin{center}
  \includegraphics[width=.38\linewidth]{question-5-1.jpg}
\end{center}
\begin{choices}
  \item $-15$
  \item $-9$
  \item $1$
  \item $9$
\end{choices}
\begin{solution}
\textbf{答案：}$A$

\textbf{分析：}画出约束条件的可行域，利用目标函数的最优解求解目标函数的最小值即可。

\textbf{详解：}$x$、$y$ 满足约束条件的可行域如图：$z=2x+y$ 经过可行域的 $A$ 时，目标函数取得最小值，由 $\begin{cases}2x-3y+3=0\\y+3=0\end{cases}$ 解得 $A(-6,-3)$，则 $z=2x+y$ 的最小值是：$-15$．故选：$A$。
\end{solution}
\end{question}

\begin{question}[points=5]
安排 $3$ 名志愿者完成 $4$ 项工作，每人至少完成 $1$ 项，每项工作由 $1$ 人完成，则不同的安排方式共有（   ）
\begin{choices}
  \item $12$ 种
  \item $18$ 种
  \item $24$ 种
  \item $36$ 种
\end{choices}
\begin{solution}
\textbf{答案：}$D$

\textbf{分析：}把工作分成 $3$ 组，然后安排工作方式即可。

\textbf{详解：}$4$ 项工作分成 $3$ 组，可得：$C_4^2=6$，安排 $3$ 名志愿者完成 $4$ 项工作，每人至少完成 $1$ 项，每项工作由 $1$ 人完成，可得：$6\times A_3^3=36$ 种．故选：$D$。
\end{solution}
\end{question}

\begin{question}[points=5]
甲、乙、丙、丁四位同学一起去问老师询问成语竞赛的成绩．老师说：你们四人中有 $2$ 位优秀，$2$ 位良好，我现在给甲看乙、丙的成绩，给乙看丙的成绩，给丁看甲的成绩．看后甲对大家说：我还是不知道我的成绩．根据以上信息，则（   ）
\begin{choices}
  \item 乙可以知道四人的成绩
  \item 丁可以知道四人的成绩
  \item 乙、丁可以知道对方的成绩
  \item 乙、丁可以知道自己的成绩
\end{choices}
\begin{solution}
\textbf{答案：}$D$

\textbf{分析：}根据四人所知只有自己看到，老师所说及最后甲说话，继而可以推出正确答案。

\textbf{详解：}四人所知只有自己看到，老师所说及最后甲说话，甲不知自己的成绩$\rightarrow$乙丙必有一优一良，（若为两优，甲会知道自己的成绩；若是两良，甲也会知道自己的成绩）$\rightarrow$乙看到了丙的成绩，知自己的成绩$\rightarrow$丁看到甲、丁也为一优一良，丁知自己的成绩，给甲看乙丙成绩，甲不知道自已的成绩，说明乙丙一优一良，假定乙丙都是优，则甲是良，假定乙丙都是良，则甲是优，那么甲就知道自已的成绩了．给乙看丙成绩，乙没有说不知道自已的成绩，假定丙是优，则乙是良，乙就知道自己成绩．给丁看甲成绩，因为甲不知道自己成绩，乙丙是一优一良，则甲丁也是一优一良，丁看到甲成绩，假定甲是优，则丁是良，丁肯定知道自已的成绩了．故选：$D$。
\end{solution}
\end{question}

\begin{question}[points=5]
执行如图的程序框图，如果输入的 $a=-1$，则输出的 $S=$（   ）

\begin{center}
  \includegraphics[width=.38\linewidth]{question-8-1.jpg}
\end{center}
\begin{choices}
  \item $2$
  \item $3$
  \item $4$
  \item $5$
\end{choices}
\begin{solution}
\textbf{答案：}$B$

\textbf{分析：}执行程序框图，依次写出每次循环得到的 $S$，$K$ 值，当 $K=7$ 时，程序终止即可得到结论。

\textbf{详解：}执行程序框图，有 $S=0$，$K=1$，$a=-1$，代入循环，第一次满足循环，$S=-1$，$a=1$，$K=2$；满足条件，第二次满足循环，$S=1$，$a=-1$，$K=3$；满足条件，第三次满足循环，$S=-2$，$a=1$，$K=4$；满足条件，第四次满足循环，$S=2$，$a=-1$，$K=5$；满足条件，第五次满足循环，$S=-3$，$a=1$，$K=6$；满足条件，第六次满足循环，$S=3$，$a=-1$，$K=7$；$K\le 6$ 不成立，退出循环输出 $S$ 的值为 $3$．故选：$B$。
\end{solution}
\end{question}

\begin{question}[points=5]
若双曲线 $C:\frac{x^2}{a^2}-\frac{y^2}{b^2}=1$（$a>0$，$b>0$）的一条渐近线被圆 $(x-2)^2+y^2=4$ 所截得的弦长为 $2$，则 $C$ 的离心率为（   ）
\begin{choices}
  \item $2$
  \item $\sqrt{3}$
  \item $\sqrt{2}$
  \item $\frac{2\sqrt{3}}{3}$
\end{choices}
\begin{solution}
\textbf{答案：}$A$

\textbf{分析：}通过圆的圆心与双曲线的渐近线的距离，列出关系式，然后求解双曲线的离心率即可。

\textbf{详解：}双曲线 $C:\frac{x^2}{a^2}-\frac{y^2}{b^2}=1$（$a>0$，$b>0$）的一条渐近线不妨为：$bx+ay=0$，圆 $(x-2)^2+y^2=4$ 的圆心 $(2,0)$，半径为：$2$，双曲线的一条渐近线被圆所截得的弦长为 $2$，可得圆心到直线的距离为：$d=\frac{|2b|}{\sqrt{a^2+b^2}}=\sqrt{2^2-1^2}=\sqrt{3}$，解得：$\frac{4b^2}{a^2+b^2}=3$，可得 $e^2=4$，即 $e=2$．故选：$A$。
\end{solution}
\end{question}

\begin{question}[points=5]
已知直三棱柱 $ABC-A_1B_1C_1$ 中，$\angle ABC=120^\circ$，$AB=2$，$BC=CC_1=1$，则异面直线 $AB_1$ 与 $BC_1$ 所成角的余弦值为（   ）
\begin{choices}
  \item $\frac{\sqrt{15}}{5}$
  \item $\frac{\sqrt{10}}{5}$
  \item $\frac{\sqrt{10}}{5}$
  \item $\frac{\sqrt{5}}{5}$
\end{choices}
\begin{solution}
\textbf{答案：}$C$

\textbf{分析：}【解法一】设 $M$、$N$、$P$ 分别为 $AB$，$BB_1$ 和 $B_1C_1$ 的中点，得出 $AB_1$、$BC_1$ 夹角为 $MN$ 和 $NP$ 夹角或其补角；根据中位线定理，结合余弦定理求出 $AC$、$MQ$，$MP$ 和 $\angle MNP$ 的余弦值即可．【解法二】通过补形的办法，把原来的直三棱柱变成直四棱柱，解法更简洁。

\textbf{详解：}【解法一】如图所示，设 $M$、$N$、$P$ 分别为 $AB$，$BB_1$ 和 $B_1C_1$ 的中点，则 $AB_1$、$BC_1$ 夹角为 $MN$ 和 $NP$ 夹角或其补角（因异面直线所成角为 $(0,\frac{\pi}{2}]$），可知 $MN=\frac{1}{2}AB_1=\frac{\sqrt{5}}{2}$，$NP=\frac{1}{2}BC_1=\frac{\sqrt{2}}{2}$；作 $BC$ 中点 $Q$，则 $\triangle PQM$ 为直角三角形；$\because PQ=1$，$MQ=\frac{1}{2}AC$，$\triangle ABC$ 中，由余弦定理得 $AC^2=AB^2+BC^2-2AB\cdot BC\cdot\cos\angle ABC=4+1-2\times2\times1\times(-\frac{1}{2})=7$，$\therefore AC=\sqrt{7}$，$\therefore MQ=\frac{\sqrt{7}}{2}$；在 $\triangle MQP$ 中，$MP=\sqrt{PQ^2+MQ^2}=\frac{\sqrt{11}}{2}$；在 $\triangle PMN$ 中，由余弦定理得 $\cos\angle MNP=\frac{MN^2+NP^2-PM^2}{2\cdot MN\cdot NP}=\frac{\frac{5}{4}+\frac{1}{2}-\frac{11}{4}}{2\times\frac{\sqrt{5}}{2}\times\frac{\sqrt{2}}{2}}=\frac{-\frac{1}{2}}{\sqrt{10}}=-\frac{\sqrt{10}}{10}$；又异面直线所成角的范围是 $(0,\frac{\pi}{2}]$，$\therefore AB_1$ 与 $BC_1$ 所成角的余弦值为 $\frac{\sqrt{10}}{5}$．【解法二】如图所示，补成四棱柱 $ABCD-A_1B_1C_1D_1$，求 $\angle BC_1D$ 即可；$BC_1=\sqrt{2}$，$BD=\sqrt{1^2+1^2-2\times1\times1\times\cos60^\circ}=1$，$C_1D=\sqrt{5}$，$\because BC_1^2+BD^2=C_1D^2$，$\therefore\angle DBC_1=90^\circ$，$\therefore\cos\angle BC_1D=\frac{BC_1}{C_1D}=\frac{\sqrt{2}}{\sqrt{5}}=\frac{\sqrt{10}}{5}$．故选：$C$。
\end{solution}
\end{question}

\begin{question}[points=5]
若 $x=-2$ 是函数 $f(x)=(x^2+ax-1)e^{x-1}$ 的极值点，则 $f(x)$ 的极小值为（   ）
\begin{choices}
  \item $-1$
  \item $-2e^{-3}$
  \item $5e^{-3}$
  \item $1$
\end{choices}
\begin{solution}
\textbf{答案：}$A$

\textbf{分析：}求出函数的导数，利用极值点，求出 $a$，然后判断函数的单调性，求解函数的极小值即可。

\textbf{详解：}函数 $f(x)=(x^2+ax-1)e^{x-1}$，可得 $f'(x)=(2x+a)e^{x-1}+(x^2+ax-1)e^{x-1}$，$x=-2$ 是函数 $f(x)=(x^2+ax-1)e^{x-1}$ 的极值点，可得：$f'(-2)=(-4+a)e^{-3}+(4-2a-1)e^{-3}=0$，即 $-4+a+(3-2a)=0$．解得 $a=-1$．可得 $f'(x)=(2x-1)e^{x-1}+(x^2-x-1)e^{x-1}=(x^2+x-2)e^{x-1}$，函数的极值点为：$x=-2$，$x=1$，当 $x<-2$ 或 $x>1$ 时，$f'(x)>0$ 函数是增函数，$x\in(-2,1)$ 时，函数是减函数，$x=1$ 时，函数取得极小值：$f(1)=(1^2-1-1)e^{1-1}=-1$．故选：$A$。
\end{solution}
\end{question}

\begin{question}[points=5]
已知 $\triangle ABC$ 是边长为 $2$ 的等边三角形，$P$ 为平面 $ABC$ 内一点，则 $\overrightarrow{PA}\cdot(\overrightarrow{PB}+\overrightarrow{PC})$ 的最小值是（   ）
\begin{choices}
  \item $-2$
  \item $-\frac{3}{2}$
  \item $-\frac{4}{3}$
  \item $-1$
\end{choices}
\begin{solution}
\textbf{答案：}$B$

\textbf{分析：}根据条件建立坐标系，求出点的坐标，利用坐标法结合向量数量积的公式进行计算即可。

\textbf{详解：}建立如图所示的坐标系，以 $BC$ 中点为坐标原点，则 $A(0,\sqrt{3})$，$B(-1,0)$，$C(1,0)$，设 $P(x,y)$，则 $\overrightarrow{PA}=(-x,\sqrt{3}-y)$，$\overrightarrow{PB}=(-1-x,-y)$，$\overrightarrow{PC}=(1-x,-y)$，则 $\overrightarrow{PA}\cdot(\overrightarrow{PB}+\overrightarrow{PC})=2x^2-2\sqrt{3}y+2y^2=2[x^2+(y-\frac{\sqrt{3}}{2})^2-\frac{3}{4}]$，$\therefore$ 当 $x=0$，$y=\frac{\sqrt{3}}{2}$ 时，取得最小值 $2\times(-\frac{3}{4})=-\frac{3}{2}$，故选：$B$。
\end{solution}
\end{question}

\section*{二、填空题}
本题共 4 小题，每小题 5 分，共 20 分。

\begin{question}[points=5]
一批产品的二等品率为 $0.02$，从这批产品中每次随机取一件，有放回地抽取 $100$ 次．$X$ 表示抽到的二等品件数，则 $DX=$ ．
\fillin[{$1.96$}]
\begin{solution}
\textbf{答案：}$1.96$

\textbf{分析：}判断概率满足的类型，然后求解方差即可。

\textbf{详解：}由题意可知，该事件满足独立重复试验，是一个二项分布模型，其中，$p=0.02$，$n=100$，则 $DX=npq=np(1-p)=100\times0.02\times0.98=1.96$．故答案为：$1.96$。
\end{solution}
\end{question}

\begin{question}[points=5]
函数 $f(x)=\sin^2x+\sqrt{3}\cos x-\frac{3}{4}$（$x\in[0,\frac{\pi}{2}]$）的最大值是 ．
\fillin[{$1$}]
\begin{solution}
\textbf{答案：}$1$

\textbf{分析：}同角的三角函数的关系以及二次函数的性质即可求出。

\textbf{详解：}$f(x)=\sin^2x+\sqrt{3}\cos x-\frac{3}{4}=1-\cos^2x+\sqrt{3}\cos x-\frac{3}{4}$，令 $\cos x=t$ 且 $t\in[0,1]$，则 $y=-t^2+\sqrt{3}t+\frac{1}{4}=-(t-\frac{\sqrt{3}}{2})^2+1$，当 $t=\frac{\sqrt{3}}{2}$ 时，$f(t)_{\max}=1$，即 $f(x)$ 的最大值为 $1$，故答案为：$1$。
\end{solution}
\end{question}

\begin{question}[points=5]
等差数列 $\{a_n\}$ 的前 $n$ 项和为 $S_n$，$a_3=3$，$S_4=10$，则 $\sum_{k=1}^{n}\frac{1}{S_k}=$ ．
\fillin[{$\frac{2n}{n+1}$}]
\begin{solution}
\textbf{答案：}$\frac{2n}{n+1}$

\textbf{分析：}利用已知条件求出等差数列的前 $n$ 项和，然后化简所求的表达式，求解即可。

\textbf{详解：}等差数列 $\{a_n\}$ 的前 $n$ 项和为 $S_n$，$a_3=3$，$S_4=10$，$S_4=2(a_2+a_3)=10$，可得 $a_2=2$，数列的首项为 $1$，公差为 $1$，$S_n=\frac{n(n+1)}{2}$，$\frac{1}{S_k}=\frac{2}{k(k+1)}=2(\frac{1}{k}-\frac{1}{k+1})$，则 $\sum_{k=1}^{n}\frac{1}{S_k}=2[1-\frac{1}{2}+\frac{1}{2}-\frac{1}{3}+\cdots+\frac{1}{n}-\frac{1}{n+1}]=2(1-\frac{1}{n+1})=\frac{2n}{n+1}$．故答案为：$\frac{2n}{n+1}$。
\end{solution}
\end{question}

\begin{question}[points=5]
已知 $F$ 是抛物线 $C:y^2=8x$ 的焦点，$M$ 是 $C$ 上一点，$FM$ 的延长线交 $y$ 轴于点 $N$．若 $M$ 为 $FN$ 的中点，则 $|FN|=$ ．
\fillin[{$6$}]
\begin{solution}
\textbf{答案：}$6$

\textbf{分析：}求出抛物线的焦点坐标，推出 $M$ 坐标，然后求解即可。

\textbf{详解：}抛物线 $C:y^2=8x$ 的焦点 $F(2,0)$，$M$ 是 $C$ 上一点，$FM$ 的延长线交 $y$ 轴于点 $N$．若 $M$ 为 $FN$ 的中点，可知 $M$ 的横坐标为：$1$，则 $M$ 的纵坐标为：$\pm2\sqrt{2}$，$|FN|=2|FM|=2\sqrt{(1-2)^2+(\pm2\sqrt{2})^2}=6$．故答案为：$6$。
\end{solution}
\end{question}

\section*{三、解答题}
解答应写出文字说明、证明过程或演算步骤。第17～21题为必考题，每个试题考生都必须作答。第22、23题为选考题，考生根据要求作答。

\begin{question}[points=12]
$\triangle ABC$ 的内角 $A$，$B$，$C$ 的对边分别为 $a$，$b$，$c$，已知 $\sin(A+C)=8\sin^2\frac{B}{2}$．

（1）求 $\cos B$；

（2）若 $a+c=6$，$\triangle ABC$ 的面积为 $2$，求 $b$．
\begin{solution}
\textbf{答案：}（1）$\cos B=\frac{15}{17}$；（2）$b=2$

\textbf{分析：}（1）利用三角形的内角和定理可知 $A+C=\pi-B$，再利用诱导公式化简 $\sin(A+C)$，利用降幂公式化简 $8\sin^2\frac{B}{2}$，结合 $\sin^2B+\cos^2B=1$，求出 $\cos B$，（2）由（1）可知 $\sin B=\frac{8}{17}$，利用勾面积公式求出 $ac$，再利用余弦定理即可求出 $b$。

\textbf{详解：}（1）$\sin(A+C)=8\sin^2\frac{B}{2}$，$\therefore\sin B=4(1-\cos B)$，$\because\sin^2B+\cos^2B=1$，$\therefore16(1-\cos B)^2+\cos^2B=1$，$\therefore16(1-\cos B)^2+\cos^2B-1=0$，$\therefore16(\cos B-1)^2+(\cos B-1)(\cos B+1)=0$，$\therefore(17\cos B-15)(\cos B-1)=0$，$\therefore\cos B=\frac{15}{17}$；
（2）由（1）可知 $\sin B=\frac{8}{17}$，$\because S_{\triangle ABC}=\frac{1}{2}ac\cdot\sin B=2$，$\therefore ac=\frac{17}{2}$，$\therefore b^2=a^2+c^2-2ac\cos B=a^2+c^2-2\times\frac{17}{2}\times\frac{15}{17}=a^2+c^2-15=(a+c)^2-2ac-15=36-17-15=4$，$\therefore b=2$。
\end{solution}
\end{question}

\begin{question}[points=12]
海水养殖场进行某水产品的新、旧网箱养殖方法的产量对比，收获时各随机抽取了 $100$ 个网箱，测量各箱水产品的产量（单位：kg），其频率分布直方图如图：

（1）设两种养殖方法的箱产量相互独立，记 $A$ 表示事件“旧养殖法的箱产量低于 $50$kg，新养殖法的箱产量不低于 $50$kg”，估计 $A$ 的概率；

（2）填写下面列联表，并根据列联表判断是否有 $99\%$ 的把握认为箱产量与养殖方法有关：

$\begin{array}{c|cc} & \text{箱产量}<50\text{kg} & \text{箱产量}\ge50\text{kg} \\ \hline \text{旧养殖法} & & \\ \text{新养殖法} & & \end{array}$

（3）根据箱产量的频率分布直方图，求新养殖法箱产量的中位数的估计值（精确到 $0.01$）．

附：$P(K^2\ge k)$ 的临界值表：$k_{0.050}=3.841$，$k_{0.010}=6.635$，$k_{0.001}=10.828$，$K^2=\frac{n(ad-bc)^2}{(a+b)(c+d)(a+c)(b+d)}$．

\begin{center}
  \includegraphics[width=.38\linewidth]{question-18-1.jpg}
\end{center}
\begin{solution}
\textbf{答案：}（1）$0.4092$；（2）有 $99\%$ 的把握认为箱产量与养殖方法有关；（3）$52.35$kg

\textbf{分析：}（1）由题意可知：$P(A)=P(BC)=P(B)P(C)$，分布求得发生的频率，即可求得其概率；（2）完成 $2\times2$ 列联表：求得观测值，与参考值比较，即可求得有 $99\%$ 的把握认为箱产量与养殖方法有关：（3）根据频率分布直方图即可求得其中位数。

\textbf{详解：}（1）记 $B$ 表示事件“旧养殖法的箱产量低于 $50$kg”，$C$ 表示事件“新养殖法的箱产量不低于 $50$kg”，由 $P(A)=P(BC)=P(B)P(C)$，则旧养殖法的箱产量低于 $50$kg：$(0.012+0.014+0.024+0.034+0.040)\times5=0.62$，故 $P(B)$ 的估计值 $0.62$，新养殖法的箱产量不低于 $50$kg：$(0.068+0.046+0.010+0.008)\times5=0.66$，故 $P(C)$ 的估计值为 $0.66$，则事件 $A$ 的概率估计值为 $P(A)=P(B)P(C)=0.62\times0.66=0.4092$；
（2）$2\times2$ 列联表：
$\begin{array}{c|cc|c} & \text{箱产量}<50\text{kg} & \text{箱产量}\ge50\text{kg} & \text{总计} \\ \hline \text{旧养殖法} & 62 & 38 & 100 \\ \text{新养殖法} & 34 & 66 & 100 \\ \hline \text{总计} & 96 & 104 & 200 \end{array}$
则 $K^2=\frac{200\times(62\times66-38\times34)^2}{100\times100\times96\times104}\approx15.705$，由 $15.705>6.635$，$\therefore$ 有 $99\%$ 的把握认为箱产量与养殖方法有关；
（3）由新养殖法的箱产量频率分布直方图中，箱产量低于 $50$kg 的直方图的面积：$(0.004+0.020+0.044)\times5=0.34$，箱产量低于 $55$kg 的直方图面积为：$(0.004+0.020+0.044+0.068)\times5=0.68>0.5$，故新养殖法产量的中位数的估计值为：$50+\frac{0.5-0.34}{0.068}\approx52.35$（kg），新养殖法箱产量的中位数的估计值 $52.35$（kg）。
\end{solution}
\end{question}

\begin{question}[points=12]
如图，四棱锥 $P-ABCD$ 中，侧面 $PAD$ 为等边三角形且垂直于底面 $ABCD$，$AB=BC=\frac{1}{2}AD$，$\angle BAD=\angle ABC=90^\circ$，$E$ 是 $PD$ 的中点．

（1）证明：直线 $CE\parallel$ 平面 $PAB$；

（2）点 $M$ 在棱 $PC$ 上，且直线 $BM$ 与底面 $ABCD$ 所成角为 $45^\circ$，求二面角 $M-AB-D$ 的余弦值．

\begin{center}
  \includegraphics[width=.38\linewidth]{question-19-1.jpg}
\end{center}
\begin{solution}
\textbf{答案：}（1）略；（2）$\frac{\sqrt{10}}{5}$

\textbf{分析：}（1）取 $PA$ 的中点 $F$，连接 $EF$，$BF$，通过证明 $CE\parallel BF$，利用直线与平面平行的判定定理证明即可．（2）利用已知条件转化求解 $M$ 到底面的距离，作出二面角的平面角，然后求解二面角 $M-AB-D$ 的余弦值即可。

\textbf{详解：}（1）证明：取 $PA$ 的中点 $F$，连接 $EF$，$BF$，因为 $E$ 是 $PD$ 的中点，所以 $EF\parallel\frac{1}{2}AD$，$AB=BC=\frac{1}{2}AD$，$\angle BAD=\angle ABC=90^\circ$，$\therefore BC\parallel\frac{1}{2}AD$，$\therefore BCEF$ 是平行四边形，可得 $CE\parallel BF$，$BF\subset$ 平面 $PAB$，$CE\not\subset$ 平面 $PAB$，$\therefore$ 直线 $CE\parallel$ 平面 $PAB$；
（2）解：四棱锥 $P-ABCD$ 中，侧面 $PAD$ 为等边三角形且垂直于底面 $ABCD$，$AB=BC=\frac{1}{2}AD$，$\angle BAD=\angle ABC=90^\circ$，$E$ 是 $PD$ 的中点．取 $AD$ 的中点 $O$，$M$ 在底面 $ABCD$ 上的射影 $N$ 在 $OC$ 上，设 $AD=2$，则 $AB=BC=1$，$OP=\sqrt{3}$，$\therefore\angle PCO=60^\circ$，直线 $BM$ 与底面 $ABCD$ 所成角为 $45^\circ$，可得：$BN=MN$，$CN=\frac{\sqrt{3}}{3}MN$，$BC=1$，可得：$1+\frac{1}{3}BN^2=BN^2$，$BN=\frac{\sqrt{6}}{2}$，$MN=\frac{\sqrt{6}}{2}$，作 $NQ\perp AB$ 于 $Q$，连接 $MQ$，$AB\perp MN$，所以 $\angle MQN$ 就是二面角 $M-AB-D$ 的平面角，$MQ=\sqrt{MN^2+NQ^2}=\sqrt{(\frac{\sqrt{6}}{2})^2+(\frac{1}{2})^2}=\frac{\sqrt{7}}{2}$，二面角 $M-AB-D$ 的余弦值为：$\frac{NQ}{MQ}=\frac{\frac{1}{2}}{\frac{\sqrt{7}}{2}}=\frac{\sqrt{7}}{7}$。
\end{solution}
\end{question}

\begin{question}[points=12]
设 $O$ 为坐标原点，动点 $M$ 在椭圆 $C:\frac{x^2}{2}+y^2=1$ 上，过 $M$ 作 $x$ 轴的垂线，垂足为 $N$，点 $P$ 满足 $\overrightarrow{NP}=\sqrt{2}\overrightarrow{NM}$．

（1）求点 $P$ 的轨迹方程；

（2）设点 $Q$ 在直线 $x=-3$ 上，且 $\overrightarrow{OP}\cdot\overrightarrow{PQ}=1$．证明：过点 $P$ 且垂直于 $OQ$ 的直线 $l$ 过 $C$ 的左焦点 $F$．
\begin{solution}
\textbf{答案：}（1）$x^2+y^2=2$；（2）略

\textbf{分析：}（1）设 $M(x_0,y_0)$，由题意可得 $N(x_0,0)$，设 $P(x,y)$，运用向量的坐标运算，结合 $M$ 满足椭圆方程，化简整理可得 $P$ 的轨迹方程；（2）设 $Q(-3,m)$，$P(\sqrt{2}\cos\alpha,\sin\alpha)$，$0\le\alpha<2\pi$，运用向量的数量积的坐标表示，可得 $m$，即有 $Q$ 的坐标，求得椭圆的左焦点坐标，求得 $OQ$，$PF$ 的斜率，由两直线垂直的条件：向量数量积为 $0$，即可得证。

\textbf{详解：}（1）设 $M(x_0,y_0)$，由题意可得 $N(x_0,0)$，设 $P(x,y)$，由点 $P$ 满足 $\overrightarrow{NP}=\sqrt{2}\overrightarrow{NM}$．可得 $(x-x_0,y)=\sqrt{2}(0,y_0)$，可得 $x-x_0=0$，$y=\sqrt{2}y_0$，即有 $x_0=x$，$y_0=\frac{y}{\sqrt{2}}$，代入椭圆方程 $\frac{x^2}{2}+y^2=1$，可得 $\frac{x^2}{2}+\frac{y^2}{2}=1$，即有点 $P$ 的轨迹方程为圆 $x^2+y^2=2$；
（2）证明：设 $Q(-3,m)$，$P(\sqrt{2}\cos\alpha,\sin\alpha)$（$0\le\alpha<2\pi$），$\overrightarrow{OP}\cdot\overrightarrow{PQ}=1$，可得 $(\sqrt{2}\cos\alpha,\sin\alpha)\cdot(-3-\sqrt{2}\cos\alpha,m-\sin\alpha)=1$，即为 $-3\sqrt{2}\cos\alpha-2\cos^2\alpha+\sqrt{2}m\sin\alpha-2\sin^2\alpha=1$，当 $\alpha=0$ 时，上式不成立，则 $0<\alpha<2\pi$，解得 $m=\frac{3(1+\sqrt{2}\cos\alpha)}{\sqrt{2}\sin\alpha}$，即有 $Q(-3,\frac{3(1+\sqrt{2}\cos\alpha)}{\sqrt{2}\sin\alpha})$，椭圆 $\frac{x^2}{2}+y^2=1$ 的左焦点 $F(-1,0)$，由 $\overrightarrow{PF}\cdot\overrightarrow{OQ}=(-1-\sqrt{2}\cos\alpha,-\sin\alpha)\cdot(-3,\frac{3(1+\sqrt{2}\cos\alpha)}{\sqrt{2}\sin\alpha})=3+3\sqrt{2}\cos\alpha-3(1+\sqrt{2}\cos\alpha)=0$．可得过点 $P$ 且垂直于 $OQ$ 的直线 $l$ 过 $C$ 的左焦点 $F$．
\end{solution}
\end{question}

\begin{question}[points=12]
已知函数 $f(x)=ax^2-ax-x\ln x$，且 $f(x)\ge0$．

（1）求 $a$；

（2）证明：$f(x)$ 存在唯一的极大值点 $x_0$，且 $e^{-2}<f(x_0)<2^{-2}$．
\begin{solution}
\textbf{答案：}（1）$a=1$；（2）略

\textbf{分析：}（1）通过分析可知 $f(x)\ge0$ 等价于 $h(x)=ax-a-\ln x\ge0$，进而利用 $h'(x)=a-\frac{1}{x}$ 可得 $h(x)_{\min}=h(\frac{1}{a})$，从而可得结论；（2）通过（1）可知 $f(x)=x^2-x-x\ln x$，记 $t(x)=f'(x)=2x-2-\ln x$，解不等式可知 $t(x)_{\min}=t(\frac{1}{2})=\ln2-1<0$，从而可知 $f'(x)=0$ 存在两根 $x_0$，$x_2$，利用 $f(x)$ 必存在唯一极大值点 $x_0$ 及 $x_0<\frac{1}{2}$ 可知 $f(x_0)<\frac{1}{4}$，另一方面可知 $f(x_0)>f(\frac{1}{e})=\frac{1}{e^2}$。

\textbf{详解：}（1）解：因为 $f(x)=ax^2-ax-x\ln x=x(ax-a-\ln x)$（$x>0$），则 $f(x)\ge0$ 等价于 $h(x)=ax-a-\ln x\ge0$，求导可知 $h'(x)=a-\frac{1}{x}$．则当 $a\le0$ 时 $h'(x)<0$，即 $y=h(x)$ 在 $(0,+\infty)$ 上单调递减，所以当 $x_0>1$ 时，$h(x_0)<h(1)=0$，矛盾，故 $a>0$．因为当 $0<x<\frac{1}{a}$ 时 $h'(x)<0$、当 $x>\frac{1}{a}$ 时 $h'(x)>0$，所以 $h(x)_{\min}=h(\frac{1}{a})$，又因为 $h(1)=a-a-\ln1=0$，所以 $\frac{1}{a}=1$，解得 $a=1$；
（2）证明：由（1）可知 $f(x)=x^2-x-x\ln x$，$f'(x)=2x-2-\ln x$，令 $f'(x)=0$，可得 $2x-2-\ln x=0$，记 $t(x)=2x-2-\ln x$，则 $t'(x)=2-\frac{1}{x}$，令 $t'(x)=0$，解得：$x=\frac{1}{2}$，所以 $t(x)$ 在区间 $(0,\frac{1}{2})$ 上单调递减，在 $(\frac{1}{2},+\infty)$ 上单调递增，所以 $t(x)_{\min}=t(\frac{1}{2})=\ln2-1<0$，从而 $t(x)=0$ 有解，即 $f'(x)=0$ 存在两根 $x_0$，$x_2$，且不妨设 $f'(x)$ 在 $(0,x_0)$ 上为正、在 $(x_0,x_2)$ 上为负、在 $(x_2,+\infty)$ 上为正，所以 $f(x)$ 必存在唯一极大值点 $x_0$，且 $2x_0-2-\ln x_0=0$，所以 $f(x_0)=x_0^2-x_0-x_0\ln x_0=x_0^2-x_0+2x_0-2x_0^2=x_0-x_0^2$，由 $x_0<\frac{1}{2}$ 可知 $f(x_0)<(x_0-x_0^2)_{\max}=-\frac{1}{4}+\frac{1}{2}=\frac{1}{4}$；由 $f'(\frac{1}{e})<0$ 可知 $x_0<\frac{1}{e}<\frac{1}{2}$，所以 $f(x)$ 在 $(0,x_0)$ 上单调递增，在 $(x_0,\frac{1}{e})$ 上单调递减，所以 $f(x_0)>f(\frac{1}{e})=\frac{1}{e^2}$；综上所述，$f(x)$ 存在唯一的极大值点 $x_0$，且 $e^{-2}<f(x_0)<2^{-2}$。
\end{solution}
\end{question}

\begin{question}[points=10]
[选修4-4：坐标系与参数方程]

在直角坐标系 $xOy$ 中，以坐标原点为极点，$x$ 轴的正半轴为极轴建立极坐标系，曲线 $C_1$ 的极坐标方程为 $\rho\cos\theta=4$．

（1）$M$ 为曲线 $C_1$ 上的动点，点 $P$ 在线段 $OM$ 上，且满足 $|OM|\cdot|OP|=16$，求点 $P$ 的轨迹 $C_2$ 的直角坐标方程；

（2）设点 $A$ 的极坐标为 $(2,\frac{\pi}{3})$，点 $B$ 在曲线 $C_2$ 上，求 $\triangle OAB$ 面积的最大值．
\begin{solution}
\textbf{答案：}（1）$(x-2)^2+y^2=4$（$x\neq0$）；（2）$2+\sqrt{3}$

\textbf{分析：}（1）设 $P(x,y)$，利用相似得出 $M$ 点坐标，根据 $|OM|\cdot|OP|=16$ 列方程化简即可；（2）求出曲线 $C_2$ 的圆心和半径，得出 $B$ 到 $OA$ 的最大距离，即可得出最大面积。

\textbf{详解：}（1）曲线 $C_1$ 的直角坐标方程为：$x=4$，设 $P(x,y)$，$M(4,y_0)$，则 $\frac{y}{y_0}=\frac{x}{4}$，$\therefore y_0=\frac{4y}{x}$，$\because|OM||OP|=16$，$\therefore\sqrt{16+y_0^2}\cdot\sqrt{x^2+y^2}=16$，即 $(x^2+y^2)(1+\frac{16y^2}{x^2})=16$，$\therefore x^4+2x^2y^2+y^4=16x^2$，即 $(x^2+y^2)^2=16x^2$，两边开方得：$x^2+y^2=4x$，整理得：$(x-2)^2+y^2=4$（$x\neq0$），$\therefore$ 点 $P$ 的轨迹 $C_2$ 的直角坐标方程：$(x-2)^2+y^2=4$（$x\neq0$）；
（2）点 $A$ 的直角坐标为 $A(1,\sqrt{3})$，显然点 $A$ 在曲线 $C_2$ 上，$|OA|=2$，$\therefore$ 曲线 $C_2$ 的圆心 $(2,0)$ 到弦 $OA$ 的距离 $d=\sqrt{2^2-(\frac{2}{2})^2}=\sqrt{3}$，$\therefore\triangle AOB$ 的最大面积 $S=\frac{1}{2}|OA|\cdot(2+\sqrt{3})=2+\sqrt{3}$。
\end{solution}
\end{question}

\begin{question}[points=10]
[选修4-5：不等式选讲]

已知 $a>0$，$b>0$，$a^3+b^3=2$．证明：

（1）$(a+b)(a^5+b^5)\ge4$；

（2）$a+b\le2$．
\begin{solution}
\textbf{答案：}略

\textbf{分析：}（1）由柯西不等式即可证明，（2）由 $a^3+b^3=2$ 转化为 $\frac{(a+b)^3-2}{3(a+b)}=ab$，再由均值不等式可得：$\frac{(a+b)^3-2}{3(a+b)}=ab\le(\frac{a+b}{2})^2$，即可得到 $\frac{1}{4}(a+b)^3\le2$，问题得以证明。

\textbf{详解：}（1）由柯西不等式得：$(a+b)(a^5+b^5)\ge(\sqrt{a\cdot a^5}+\sqrt{b\cdot b^5})^2=(a^3+b^3)^2\ge4$，当且仅当 $\frac{a}{a^5}=\frac{b}{b^5}$，即 $a=b=1$ 时取等号；
（2）$\because a^3+b^3=2$，$\therefore(a+b)(a^2-ab+b^2)=2$，$\therefore(a+b)[(a+b)^2-3ab]=2$，$\therefore(a+b)^3-3ab(a+b)=2$，$\therefore\frac{(a+b)^3-2}{3(a+b)}=ab$，由均值不等式可得：$\frac{(a+b)^3-2}{3(a+b)}=ab\le(\frac{a+b}{2})^2$，$\therefore(a+b)^3-2\le\frac{3}{4}(a+b)^3$，$\therefore\frac{1}{4}(a+b)^3\le2$，$\therefore a+b\le2$，当且仅当 $a=b=1$ 时等号成立。
\end{solution}
\end{question}

